%% FullCourse_10Lecs -- Lecture 10, Topic 10.2: Case Study 1 --- Paper Review Skill
%% Source: ./archive_pre_split/Lec10_Case_Studies_v2.tex (section 2)

%% Shared preamble for the FullCourse_10Lecs topic decks.
%%
%% Each topic .tex \input's this file, then defines its own \title, \date
%% override (if any), \graphicspath, and \begin{document}.
%%
%% Reconstructed verbatim from the inlined copy preserved in
%% Lec10_Case_Studies/Lec10_T8_Case6_DeepRamsey.tex (lines 23-190), which is
%% explicitly annotated there as "from ../shared_preamble.tex, copied verbatim".
%% Base preamble matches MiniCourse_8hr/Slides/shared_preamble.tex; this version
%% additionally provides \sectiondivider, the conceptbox/cautionbox/examplebox/
%% takeawaybox tcolorboxes, and \compacttables used across the split decks.

\documentclass[10pt,english,aspectratio=169]{beamer}

\usepackage{ae,aecompl}
\usepackage[T1]{fontenc}
\usepackage{textcomp}
\usepackage[utf8]{inputenc}
\usepackage{babel}
\usepackage{datetime2}

\setcounter{secnumdepth}{3}
\setcounter{tocdepth}{3}

\usepackage{amsmath, amssymb, amsthm, graphicx}
\usepackage{booktabs, multirow}
\usepackage{tikz}
\usepackage{pgfplots}
\pgfplotsset{compat=1.16}
\usepackage[most]{tcolorbox}
\tcbuselibrary{skins, breakable, theorems}
\usepackage{listings}
\usepackage{xcolor}

\lstset{
  basicstyle=\ttfamily\small,
  keywordstyle=\color{blue!80!black}\bfseries,
  commentstyle=\color{green!50!black}\itshape,
  stringstyle=\color{orange!80!black},
  backgroundcolor=\color{gray!8},
  frame=single,
  rulecolor=\color{gray!40},
  breaklines=true,
  showstringspaces=false,
  numbers=none,
}

\lstdefinestyle{pythonstyle}{
    language=Python,
    basicstyle=\ttfamily\footnotesize,
    keywordstyle=\color{blue!80!black}\bfseries,
    commentstyle=\color{green!50!black}\itshape,
    stringstyle=\color{orange!80!black},
    frame=single,
    breaklines=true,
    showstringspaces=false,
    numbers=left,
    numbersep=5pt,
    numberstyle=\tiny\color{gray},
    backgroundcolor=\color{gray!8},
    rulecolor=\color{gray!40}
}

\ifx\hypersetup\undefined
  \AtBeginDocument{%
    \hypersetup{bookmarks=false,breaklinks=true,pdfborder={0 0 0},colorlinks=true,
      linkcolor=DarkText,urlcolor=RedTitle}%
  }
\else
  \hypersetup{bookmarks=false,breaklinks=true,pdfborder={0 0 0},colorlinks=true,
    linkcolor=DarkText,urlcolor=RedTitle}%
\fi

\makeatletter
\newcommand\makebeamertitle{%
  \begin{frame}[plain]
    \vspace*{1.0cm}
    \centering
    {\usebeamerfont{title}\usebeamercolor[fg]{title}\inserttitle\par}
    \vspace{1.0cm}
    {\usebeamerfont{author}\usebeamercolor[fg]{author}\insertauthor\par}
    \vspace{0.2cm}
    {\usebeamerfont{date}\usebeamercolor[fg]{date}\insertdate\par}
  \end{frame}}%
\AtBeginDocument{%
  \let\origtableofcontents=\tableofcontents
  \def\tableofcontents{\@ifnextchar[{\origtableofcontents}{\gobbletableofcontents}}
  \def\gobbletableofcontents#1{\origtableofcontents}%
}

\usetheme{metropolis}
\usefonttheme{professionalfonts}

\definecolor{LightGray}{RGB}{230,230,230}
\definecolor{RedTitle}{RGB}{0,0,200}
\definecolor{VeryLightGray}{RGB}{245,245,245}
\definecolor{DarkText}{RGB}{0,0,0}

\setbeamercolor{background canvas}{bg=white}
\setbeamercolor{title}{fg=RedTitle}
\setbeamercolor{author}{fg=DarkText}
\setbeamercolor{date}{fg=DarkText}
\setbeamercolor{frametitle}{bg=VeryLightGray,fg=DarkText}
\setbeamercolor{normal text}{fg=DarkText}
\setbeamerfont{title}{size=\large}

\setbeamersize{text margin left=5mm, text margin right=5mm}
\addtobeamertemplate{frametitle}{}{\vspace{-0.5em}}
\newcommand{\reducevspace}{\vspace{-0.5cm}}

% Shared section divider and box styles for split topic decks.
\newcommand{\sectiondivider}[2]{%
\begin{frame}[plain,noframenumbering]
\begin{center}
\vfill
{\Large\textbf{#1}\par}
\vspace{0.35cm}
{\normalsize #2\par}
\vfill
\end{center}
\end{frame}
}

\newtcolorbox{conceptbox}[1][]{
  colback=blue!3,
  colframe=RedTitle!55,
  boxrule=0.35mm,
  arc=1mm,
  left=2mm,
  right=2mm,
  top=1mm,
  bottom=1mm,
  #1
}
\newtcolorbox{cautionbox}[1][]{
  colback=red!3,
  colframe=red!50!black,
  boxrule=0.35mm,
  arc=1mm,
  left=2mm,
  right=2mm,
  top=1mm,
  bottom=1mm,
  #1
}
\newtcolorbox{examplebox}[1][]{
  colback=green!3,
  colframe=green!45!black,
  boxrule=0.35mm,
  arc=1mm,
  left=2mm,
  right=2mm,
  top=1mm,
  bottom=1mm,
  #1
}
\newtcolorbox{takeawaybox}[1][]{
  colback=VeryLightGray,
  colframe=RedTitle!55,
  boxrule=0.35mm,
  arc=1mm,
  left=2mm,
  right=2mm,
  top=1mm,
  bottom=1mm,
  #1
}
\newcommand{\compacttables}{\renewcommand{\arraystretch}{1.15}}

\setbeamertemplate{navigation symbols}{}
\setbeamertemplate{footline}{%
  \hfill%
  \usebeamercolor[fg]{page number in head/foot}%
  \usebeamerfont{page number in head/foot}%
  \insertframenumber\,/\,\inserttotalframenumber\kern1em\vskip2pt%
}

\makeatother

\author{Zhigang Feng}
% "date with time": compile timestamp so each build is identifiable.
\date{\today\ \DTMcurrenttime}


% Suppress redundant auto-generated metropolis section page; \sectiondivider frames are the dividers we keep.
\metroset{sectionpage=none}

\graphicspath{{./pic/}{./figures/}{./Exercise10_ES_Fellows/plots/}{../../MiniCourse_8hr/Slides/Module4_Agentic_AI/pic/}{../}}

\usetikzlibrary{positioning,arrows.meta,shapes,calc,fit,backgrounds}

\lstdefinestyle{markdownstyle}{
    basicstyle=\ttfamily\footnotesize,
    frame=single,
    rulecolor=\color{gray!40},
    backgroundcolor=\color{gray!8},
    breaklines=true,
    showstringspaces=false,
    numbers=none,
}

\title[AI for Econ Research]{AI for Economic Research: Dynamic Models, Language, and Agents\\
Lecture 10.2: Case Study 1 --- Paper Review Skill}

\begin{document}

\makebeamertitle

%=============================================================================
\begin{frame}{Topic 10.2 Agenda}
\begin{enumerate}
  \item \textbf{Case Study 1} --- building a paper-review skill from scratch
\end{enumerate}
\end{frame}

%=============================================================================
\section{Case Study 1: Paper Review Skill}
%===========================================================

%===========================================================

\sectiondivider{Case Study 1}{A reusable quality-control workflow for paper review}

\begin{frame}{Case Study 1: Why This Becomes a Skill}
\begin{tcolorbox}[colback=blue!3, colframe=RedTitle!60]
\textbf{Research job:} produce a structured referee-style review for a quantitative economics paper with quote-level evidence and quality gates.
\end{tcolorbox}

\vspace{0.3cm}
\begin{itemize}
    \item You may run this workflow many times across papers, coauthors, and seminars
    \item Quality matters more than speed because a wrong criticism is costly
    \item The task has internal stages with clear dependencies
    \item That makes it a good candidate for a reusable skill rather than a one-off prompt
\end{itemize}
\end{frame}

\begin{frame}{The Naive Baseline Fails Fast}
\textbf{Tempting prompt:} ``Review this paper.''

\vspace{0.25cm}
\textbf{What usually comes back:}
\begin{itemize}
    \item generic praise and generic weaknesses
    \item no verbatim anchor to the paper
    \item no proof-checking discipline
    \item no calibration to quantitative macro standards
    \item no filter removing low-confidence comments
\end{itemize}

\vspace{0.3cm}
\textbf{The failure is not just model quality.} It is the absence of structure, specialization, and a final editorial gate.
\end{frame}

\begin{frame}[fragile]{What a High-Quality Output Should Look Like}
\begin{tcolorbox}[colback=gray!5, colframe=gray!40, title={\footnotesize Shortened review template}, fonttitle=\bfseries\footnotesize]
\begin{lstlisting}[style=markdownstyle, basicstyle=\ttfamily\scriptsize]
## Overall Feedback
1. Estimand is not fully pinned down in Section 2.3
2. Proof of Theorem 2 needs one missing justification
3. Numerical section needs a reproducibility note

## Detailed Comments
### 1. Proof step in Theorem 2
Quote: "By dominated convergence ..."
Feedback: This step needs an explicit domination argument.
\end{lstlisting}
\end{tcolorbox}

\vspace{0.25cm}
\textbf{Target properties:} issue prioritization, exact anchoring, specific fixes, and a recommendation you can defend.
\end{frame}

\begin{frame}{The 7-Phase Pipeline}
\begin{center}
\begin{tikzpicture}[
    phase/.style={rectangle, rounded corners=3pt, draw=RedTitle!70, fill=blue!5,
        minimum width=1.5cm, minimum height=0.85cm, font=\scriptsize\bfseries,
        text width=1.45cm, align=center, text=DarkText},
    arr/.style={-{Stealth[length=2.5mm]}, thick, gray!60},
]
\node[phase] (p1) {1\\Extract};
\node[phase, right=0.28cm of p1] (p2) {2\\Calibrate};
\node[phase, right=0.28cm of p2] (p3) {3\\Overview};
\node[phase, right=0.28cm of p3] (p4) {4\\Section\\Reviews};
\node[phase, right=0.28cm of p4] (p5) {5\\Cross-\\Section};
\node[phase, right=0.28cm of p5] (p6) {6\\Editorial\\Filter};
\node[phase, right=0.28cm of p6] (p7) {7\\Output};

\draw[arr] (p1) -- (p2);
\draw[arr] (p2) -- (p3);
\draw[arr] (p3) -- (p4);
\draw[arr] (p4) -- (p5);
\draw[arr] (p5) -- (p6);
\draw[arr] (p6) -- (p7);
\end{tikzpicture}
\end{center}

\vspace{0.25cm}
\begin{itemize}
    \item Parallelism happens inside phases, not across them
    \item The output of one phase becomes the acceptance criteria for the next
    \item This is a quality-control pipeline, not just a prompt chain
\end{itemize}
\end{frame}

\begin{frame}{Why Decomposition Matters Here}
\begin{columns}[T]
\begin{column}{0.32\textwidth}
\textbf{Context}
\begin{itemize}
    \item long paper
    \item review rules
    \item output template
\end{itemize}
One call is too coarse.
\end{column}
\begin{column}{0.32\textwidth}
\textbf{Specialization}
\begin{itemize}
    \item proof review
    \item methodology review
    \item literature fit
\end{itemize}
One prompt does not do all well.
\end{column}
\begin{column}{0.32\textwidth}
\textbf{Quality}
\begin{itemize}
    \item generic comments
    \item invented quotes
    \item duplicated issues
\end{itemize}
You need a final filter.
\end{column}
\end{columns}

\vspace{0.3cm}
\textbf{Economics translation:} this is like assigning specialized RA tasks, then forcing everything through one senior review before it goes out.
\end{frame}

\begin{frame}{The Highest-Value Guardrails}
\begin{enumerate}
    \item \textbf{Confidence gate.} Do not assert an error unless the model can derive it step by step.
    \item \textbf{Quote verification.} Comments that cite the paper must be anchored to real text.
    \item \textbf{Editorial filter.} Remove generic advice, duplicates, and stylistic fluff.
    \item \textbf{Humanization.} The final report should read like a serious economics referee report, not model boilerplate.
\end{enumerate}

\vspace{0.3cm}
\textbf{The lesson is broader than paper review:} when the stakes are high, validation must be a named stage in the workflow.
\end{frame}

\begin{frame}{Worked Example: Why the MPE Paper Is a Good Stress Test}
\begin{tcolorbox}[colback=blue!3, colframe=RedTitle!60]
\textbf{Demo paper:} \textit{Markov Equilibria with Quasi-Hyperbolic Discounting}
\end{tcolorbox}

\vspace{0.25cm}
\begin{itemize}
    \item it mixes proofs, numerical work, and dynamic economic interpretation
    \item it forces the workflow to separate theorem checking from broader conceptual review
    \item it is hard enough that a one-shot prompt should fail
\end{itemize}

\vspace{0.3cm}
\textbf{This is why case studies matter:} you only learn the workflow once you see where the naive baseline actually breaks.
\end{frame}

\begin{frame}{What the Economist Still Owns in Case Study 1}
\begin{columns}[T]
\begin{column}{0.48\textwidth}
\textbf{Agent side}
\begin{itemize}
    \item parse sections
    \item fan out specialized prompts
    \item draft comments
    \item enforce formatting
\end{itemize}
\end{column}
\begin{column}{0.48\textwidth}
\textbf{Economist side}
\begin{itemize}
    \item decide what counts as a serious objection
    \item reject false positives
    \item judge novelty and importance
    \item decide what feedback is worth sending
\end{itemize}
\end{column}
\end{columns}

\vspace{0.3cm}
\textbf{Skill logic:} you are not outsourcing judgment. You are outsourcing disciplined first-pass execution.
\end{frame}

\begin{frame}{When Should a Workflow Become a Skill?}
\small
\begin{tabular}{p{5.7cm}p{6.0cm}}
\toprule
\textbf{Build a skill when ...} & \textbf{Stay ad hoc when ...} \\
\midrule
you will repeat the workflow many times & the task is one-off \\
the workflow has named validation stages & you are still exploring what the workflow should be \\
multiple users need consistent output & only one user needs the result \\
quality failures are expensive & output format is flexible \\
\bottomrule
\end{tabular}

\vspace{0.3cm}
\textbf{Case 1 verdict:} this is a skill because repeatability and quality control are the point.
\end{frame}


\end{document}
